\documentclass[../../main.tex]{subfiles}

\begin{document}

\section{有界线性算子}\label{section:ch13-有界线性算子}

\begin{definition}[线性算子]\label{definition:ch13-线性算子}
令 $X$ 和 $Y$ 是线性空间. 映射 $T:X\to Y$ 称为线性的, 若对每对 $u,v\in X$, 以及实数 $\alpha$ 与 $\beta$,
\begin{align*}
T(\alpha u+\beta v)=\alpha T(u)+\beta T(v).
\end{align*}
\end{definition}

线性映射常称为线性算子或线性变换. 在线性代数中人们研究有限维线性空间之间的线性算子, 在定义域与值域的基选定之后, 这些算子由矩阵乘法给出. 在我们对 $1\leqslant p<\infty$ 研究 $L^{p}(E)$ 空间时, 我们考虑从 $L^{p}$ 到 $\mathbf{R}$ 的连续线性算子. 我们称这些算子为泛函并且证明刻画它们的 Riesz 表示定理.

\begin{definition}[有界算子]\label{definition:ch13-有界算子}
令 $X$ 与 $Y$ 为赋范线性空间. 线性算子 $T:X\to Y$ 称为是有界的, 若存在常数 $M\geqslant 0$ 使得
\begin{align}
\text{对所有 }u\in X,\quad\|T(u)\|\leqslant M\|u\|.\label{align:ch13-有界}
\end{align}
所有这样的 $M$ 的下确界称为 $T$ 的算子范数, 记为 $\|T\|$. 从 $X$ 到 $Y$ 的有界算子全体记为 $\mathcal{L}(X,Y)$.
\end{definition}

令 $X$ 和 $Y$ 为赋范线性空间而 $T$ 属于 $\mathcal{L}(X,Y)$. 容易看出\eqref{align:ch13-有界}对 $M=\|T\|$ 成立. 因此, 根据 $T$ 的线性,
\begin{align}
\text{对所有 }u,v\in X,\quad\|T(u)-T(v)\|\leqslant\|T\|\cdot\|u-v\|.\label{align:ch13-Lipschitz}
\end{align}
从这我们推出有界线性算子 $T$ 的连续性质:
\begin{align}
\text{若在 }X\text{ 中 }\{u_{n}\}\to u,\quad\text{则在 }Y\text{ 中 }\{T(u_{n})\}\to T(u).\label{align:ch13-连续}
\end{align}
事实上, 关于线性算子我们有以下基本结果.

\begin{theorem}\label{theorem:ch13-连续当且仅当有界}
赋范线性空间之间的线性算子是连续的当且仅当它是有界的.
\end{theorem}

\begin{proof}
令 $X$ 和 $Y$ 为赋范线性空间而 $T:X\to Y$ 是线性的. 若 $T$ 是有界的, \eqref{align:ch13-连续}告诉我们 $T$ 是连续的. 现在假定 $T:X\to Y$ 是连续的. 由于 $T$ 是线性的, $T(0)=0$. 因此, 根据在 $u=0$ 连续的 $\varepsilon-\delta$ 准则, 其中 $\varepsilon=1$, 我们可以选取 $\delta>0$, 使得若 $\|u-0\|<\delta$, 则 $\|T(u)-T(0)\|<1$, 即若 $\|u\|<\delta$, 则 $\|T(u)\|<1$. 对任何 $u\in X$, $u\neq 0$, 设 $\lambda=\delta/\|u\|$ 且观察到根据范数的正齐次性, $\|\lambda u\|\leqslant\delta$. 因此 $\|T(\lambda u)\|\leqslant 1$. 由于 $\|T(\lambda u)\|=\lambda\|T(u)\|$, 我们得出\eqref{align:ch13-有界}对 $M=1/\delta$ 成立.
\end{proof}

\begin{definition}\label{definition:ch13-算子线性组合}
令 $X$ 和 $Y$ 是线性空间. 对线性算子 $T:X\to Y$ 与 $S:X\to Y$ 以及实数 $\alpha,\beta$, 我们逐点定义 $\alpha T+\beta S:X\to Y$ 为
\begin{align}
\text{对所有 }u\in X,\quad(\alpha T+\beta S)(u)=\alpha T(u)+\beta S(u).\label{align:ch13-算子组合}
\end{align}
\end{definition}

在逐点数乘与加法下两个线性空间之间的线性算子族是线性空间.

\begin{proposition}\label{proposition:ch13-算子空间赋范}
令 $X$ 和 $Y$ 为赋范线性空间. 则从 $X$ 到 $Y$ 的有界线性算子族 $\mathcal{L}(X,Y)$ 是一个赋范线性空间.
\end{proposition}

\begin{proof}
令 $T$ 和 $S$ 属于 $\mathcal{L}(X,Y)$. 我们从 $Y$ 上的范数的三角不等式与\eqref{align:ch13-Lipschitz}推出:
\begin{align*}
\text{对所有 }u\in X,\quad\|(T+S)(u)\|\leqslant\|T(u)\|+\|S(u)\|\\
\leqslant\|T\|\|u\|+\|S\|\|u\|=(\|T\|+\|S\|)\|u\|.
\end{align*}
因此 $T+S$ 是有界的且 $\|T+S\|\leqslant\|T\|+\|S\|$. 显然对于实数 $\alpha$, $\alpha T$ 是有界的且 $\|\alpha T\|=|\alpha|\|T\|$, 而 $\|T\|=0$ 当且仅当对所有 $u\in X$, $T(u)=0$.
\end{proof}

\begin{theorem}\label{theorem:ch13-算子空间完备}
令 $X$ 和 $Y$ 为赋范线性空间. 若 $Y$ 是 Banach 空间, 则 $\mathcal{L}(X,Y)$ 也是.
\end{theorem}

\begin{proof}
令 $\{T_{n}\}$ 为 $\mathcal{L}(X,Y)$ 中的 Cauchy 序列. 令 $u$ 属于 $X$. 则由\eqref{align:ch13-Lipschitz}, 对所有指标 $n$ 和 $m$,
\begin{align*}
\|T_{n}(u)-T_{m}(u)\|=\|(T_{n}-T_{m})u\|\leqslant\|T_{n}-T_{m}\|\cdot\|u\|.
\end{align*}
因此 $\{T_{n}(u)\}$ 是 $Y$ 中的 Cauchy 序列. 根据假设, 由于 $Y$ 是完备的, 序列 $\{T_{n}(u)\}$ 收敛到 $Y$ 的某个成员, 我们把它记为 $T(u)$. 这定义了映射 $T:X\to Y$. 我们必须证明 $T$ 属于 $\mathcal{L}(X,Y)$, 且在 $\mathcal{L}(X,Y)$ 中 $\{T_{n}\}\to T$. 为证明线性, 观察到对 $X$ 中的每对 $u_{1},u_{2}$, 由于每个 $T_{n}$ 是线性的,
\begin{align*}
T(u_{1})+Tu_{2})=\lim_{n\to\infty}T_{n}(u_{1})+\lim_{n\to\infty}T_{n}(u_{2})=\lim_{n\to\infty}T_{n}(u_{1}+u_{2})=T(u_{1}+u_{2}).
\end{align*}
类似地, $T(\lambda u)=\lambda T(u)$.

我们同时证明 $T$ 的有界性以及 $\{T_{n}\}$ 在 $\mathcal{L}(X,Y)$ 中收敛到 $T$. 令 $\varepsilon>0$. 选取指标 $N$ 使得对所有 $n\geqslant N$, $k\geqslant 1$, $\|T_{n}-T_{n+k}\|<\varepsilon/2$. 因此, 根据\eqref{align:ch13-Lipschitz}, 对所有 $u\in X$,
\begin{align*}
\|T_{n}(u)-T_{n+k}(u)\|=\|(T_{n}-T_{n+k})u\|\leqslant\|T_{n}-T_{n+k}\|\cdot\|u\|<\varepsilon/2\|u\|.
\end{align*}
固定 $n\geqslant N$ 与 $u\in X$. 由于 $\lim_{k\to\infty}T_{n+k}(u)=T(u)$ 且范数是连续的, 我们得出:
\begin{align*}
\|T_{n}(u)-T(u)\|\leqslant\varepsilon/2\|u\|.
\end{align*}
特别地, 线性算子 $T_{N}-T$ 是有界的, 由于 $T_{N}$ 也是有界的, 所以 $T$ 是有界的. 此外, 对 $n\geqslant N$, $\|T_{n}-T\|<\varepsilon$. 因此在 $\mathcal{L}(X,Y)$ 中 $\{T_{n}\}\to T$.
\end{proof}

对于两个赋范线性空间 $X$ 与 $Y$, 算子 $T\in\mathcal{L}(X,Y)$ 称为同构, 若它是一对一、映上, 且有连续的逆. 对 $\mathcal{L}(X,Y)$ 中的 $T$, 若它是一对一与映上的, 它的逆是线性的. 要成为同构, 要求逆是有界的, 即逆属于 $\mathcal{L}(Y,X)$. 两个赋范线性空间是同构的, 若它们之间存在同构映射. 这是赋范线性空间之间的等价关系. 它对赋范线性空间所起的作用与同胚对拓扑空间所起的作用相同. 保持范数的同构称为等距同构: 它是一个同构而且是与范数相对应的度量结构的等距映射.

对线性算子 $T:X\to Y$, $X$ 的子空间(其中 $\{x\in X\mid T(x)=0\}$)称为 $T$ 的核, 记为 $\ker T$. 观察到 $T$ 是一对一的当且仅当 $\ker T=\{0\}$. 我们将 $T$ 的象 $T(X)$ 记为 $\operatorname{Im}T$.

\begin{exercise}\label{exercise:ch13-09}
令 $X$ 和 $Y$ 为赋范线性空间而 $T:X\to Y$ 是线性的.
\begin{enumerate}[(i)]
\item 证明 $T$ 是连续的当且仅当它在 $X$ 中的单点 $u_{0}$ 连续.
\item 证明 $T$ 是 Lipschitz 的当且仅当它是连续的.
\item 证明: 若对 $T$ 没有线性假设, (i)与(ii)都不成立.
\end{enumerate}
\end{exercise}

\begin{exercise}\label{exercise:ch13-10}
令 $X$ 和 $Y$ 为赋范线性空间而 $T\in\mathcal{L}(X,Y)$, 证明 $\|T\|$ 是映射 $T$ 的最小 Lipschitz 常数, 即最小数 $c\geqslant 0$, 使得:
\begin{align*}
\text{对所有 }u,v\in X,\quad\|T(u)-T(v)\|\leqslant c\|u-v\|.
\end{align*}
\end{exercise}

\begin{exercise}\label{exercise:ch13-11}
令 $X$ 和 $Y$ 为赋范线性空间而 $T\in\mathcal{L}(X,Y)$, 证明:
\begin{align*}
\|T\|=\sup\{\|T(u)\|\mid u\in X,\|u\|\leqslant 1\}.
\end{align*}
\end{exercise}

\begin{exercise}\label{exercise:ch13-12}
令 $X$ 和 $Y$ 为赋范线性空间, 在 $\mathcal{L}(X,Y)$ 中 $\{T_{n}\}\to T$, 且在 $X$ 中 $\{u_{n}\}\to u$. 证明在 $Y$ 中 $\{T_{n}(u_{n})\}\to T(u)$.
\end{exercise}

\begin{exercise}\label{exercise:ch13-13}
令 $X$ 为 Banach 空间, 而 $T\in\mathcal{L}(X,X)$ 满足 $\|T\|<1$.
\begin{enumerate}[(i)]
\item 用压缩映射原理证明 $I-T\in\mathcal{L}(X,X)$ 是一对一与映上的.
\item 证明 $I-T$ 是一个同构.
\end{enumerate}
\end{exercise}

\begin{exercise}\label{exercise:ch13-14}
(Neumann 级数)令 $X$ 和 $Y$ 为赋范线性空间而 $T\in\mathcal{L}(X,X)$ 满足 $\|T\|<1$. 定义 $T^{0}=\operatorname{Id}$.
\begin{enumerate}[(i)]
\item 用 $\mathcal{L}(X,X)$ 的完备性证明 $\sum_{n=0}^{\infty}T^{n}$ 在 $\mathcal{L}(X,X)$ 中收敛.
\item 证明 $(I-T)^{-1}=\sum_{n=0}^{\infty}T^{n}$.
\end{enumerate}
\end{exercise}

\begin{exercise}\label{exercise:ch13-15}
令 $X$ 和 $Y$ 为赋范线性空间而 $T\in\mathcal{L}(Y,X)$, 证明 $T$ 是同构当且仅当存在算子 $S\in\mathcal{L}(Y,X)$, 使得对每个 $u\in X$ 与 $v\in Y$,
\begin{align*}
S(T(u))=u\text{ 且 }T(S(v))=v.
\end{align*}
\end{exercise}

\begin{exercise}\label{exercise:ch13-16}
令 $X$ 和 $Y$ 为赋范线性空间而 $T\in\mathcal{L}(X,Y)$, 证明 $\ker T$ 是 $X$ 的闭子空间且 $T$ 是一对一的当且仅当 $\ker T=\{0\}$.
\end{exercise}

\begin{exercise}\label{exercise:ch13-17}
令 $(X,\rho)$ 为包含点 $x_{0}$ 的度量空间. 定义 $\operatorname{Lip}_{0}(X)$ 为 $X$ 上的在 $x_{0}$ 消失的实值 Lipschitz 函数 $f$ 的集合. 通过对 $f\in\operatorname{Lip}_{0}(X)$ 定义
\begin{align*}
\|f\|=\sup_{x\neq y}\frac{|f(x)-f(y)|}{\rho(x,y)}
\end{align*}
证明 $\operatorname{Lip}_{0}(X)$ 是赋范线性空间. 证明 $\operatorname{Lip}_{0}(X)$ 是 Banach 空间. 对每个 $x\in X$, 定义 $\operatorname{Lip}_{0}(X)$ 上的线性泛函 $F_{x}$ 为 $F_{x}(f)=f(x)$. 证明 $F_{x}$ 属于 $\mathcal{L}(\operatorname{Lip}_{0}(X),\mathbf{R})$, 且对 $x,y\in X$, $\|F_{x}-F_{y}\|=\rho(x,y)$. 因此 $X$ 与 Banach 空间 $\mathcal{L}(\operatorname{Lip}_{0}(X),\mathbf{R})$ 的子集等距. 由于完备度量空间的任何闭子集是完备的, 这给出了任何度量空间的完备化的存在的另一个证明. 它也表明了任何度量空间等距于某个赋范线性空间的子集.
\end{exercise}

\begin{exercise}\label{exercise:ch13-18}
用前一个习题证明每个赋范线性空间是 Banach 空间的稠密子空间.
\end{exercise}

\begin{exercise}\label{exercise:ch13-19}
对赋范线性空间 $X$ 以及 $T,S\in\mathcal{L}(X,X)$, 证明复合 $S\circ T$ 也属于 $\mathcal{L}(X,X)$ 且 $\|S\circ T\|\leqslant\|S\|\cdot\|T\|$.
\end{exercise}

\begin{exercise}\label{exercise:ch13-20}
令 $X$ 为赋范线性空间而 $Y$ 是 $X$ 的闭线性子空间. 证明 $\|x\|_{1}=\inf_{y\in Y}\|x-y\|$ 定义了 $X$ 上的拟范数. 由拟范数 $\|\cdot\|_{1}$(见习题 8)诱导的赋范线性空间记为 $X/Y$ 且称为 $X$ 模 $Y$ 的商空间. 证明将 $X$ 映上 $X/Y$ 的自然映射 $\varphi$ 将开集映到开集.
\end{exercise}

\begin{exercise}\label{exercise:ch13-21}
证明: 若 $X$ 是 Banach 空间而 $Y$ 是 $X$ 的闭线性真子空间, 则商空间 $X/Y$ 也是 Banach 空间, 而自然映射 $\varphi:X\to X/Y$ 的范数为 $1$.
\end{exercise}

\begin{exercise}\label{exercise:ch13-22}
令 $X$ 与 $Y$ 为赋范线性空间, $T\in\mathcal{L}(X,Y)$ 且 $\ker T=Z$. 证明存在唯一的从 $X/Z$ 到 $Y$ 的有界线性算子 $S$, 使得 $T=S\circ\varphi$, 其中 $\varphi:X\to X/Z$ 是自然映射. 此外, 证明 $\|T\|=\|S\|$.
\end{exercise}

\end{document}